\section{Introduction}
\label{sec:introduzione}

\subsection{Thermochemical non-equilibrium in atmospheric entry}
\label{sec:contesto}

A vehicle entering the atmosphere at hypersonic speed generates a strong shock
wave. The post-shock temperature can reach tens of thousands of kelvin,
leading to vibrational and electronic excitation and molecular dissociation.
These processes occur on different time scales. Translation and rotation
equilibrate within a few molecular collisions, whereas vibrational relaxation
requires many more collisions, resulting in thermal non-equilibrium.
Park's two-temperature model~\cite{park1990,park1993} represents this behaviour
using separate translational-rotational and vibro-electronic temperatures,
coupled by energy exchange governed by a relaxation time. Chemical reactions
introduce an additional coupling: vibrationally excited molecules dissociate
more readily, while reactions remove or release vibrational energy. The
representation of these processes directly affects predictions of the heat
loads on the vehicle.

\subsection{The one-temperature model and its limitations}
\label{sec:una-temperatura}

Molecular energy is distributed among translational, rotational, vibrational
and electronic modes. A conventional one-temperature CFD formulation,
including that used by the standard OpenFOAM solvers, describes all modes
using a common temperature. This formulation assumes instantaneous energy
redistribution among the modes according to their equilibrium populations.
The same temperature determines the modal energies and the chemical reaction
rates. The assumption is appropriate when collisional energy redistribution
occurs on a time scale much shorter than the characteristic flow time.

Behind a hypersonic shock, this separation of time scales may no longer hold.
Energy is initially deposited in translational motion and is transferred to
the vibrational modes through successive collisions over a time comparable to
the residence time in the post-shock region. A one-temperature model cannot
resolve this delay: it underpredicts the translational temperature immediately
behind the shock and does not represent the subsequent thermal relaxation
region. Chemical kinetics are also affected, since dissociation rates depend
on both collision energy and molecular vibrational excitation, which a single
temperature cannot distinguish. These limitations introduce errors in the
predicted gas temperature and composition, with consequences for the shock
position and wall heat flux~\cite{park1990}.

\subsection{The two-temperature model}
\label{sec:soluzione}

Park's two-temperature model~\cite{park1990,park1993} addresses these
limitations by assigning the temperature $\Ttr$ to translation and rotation
and the temperature $\Tve$ to vibration and electronic excitation. The two
temperatures can differ by thousands of kelvin immediately behind the shock.
Energy exchange between the two groups of modes progressively reduces this
difference until thermal equilibrium is reached; relaxation is generally
faster at higher temperatures and densities. Chemical kinetics depend on both
temperatures. Dissociation rates are evaluated at an intermediate temperature
to account for the reduced dissociation rate of molecules with low vibrational
excitation, and the vibrational energy associated with dissociating molecules
is included in the energy balance. The model assumes equilibrium within each
group of modes and therefore provides an approximate description of thermal
non-equilibrium without resolving individual energy levels.

\subsection{The reference study}
\label{sec:riferimento}

The reference study by Casseau
\textit{et al.}~\cite{casseau2016}, from the University of Strathclyde,
introduces \hyfoam~\cite{hystrath}, an open-source two-temperature solver for
reacting hypersonic flows based on OpenFOAM~\cite{openfoam13}. Its initial
verification uses a series of heat-bath cases, in which isolated, stationary
gas volumes evolve from prescribed non-equilibrium conditions. These cases
allow the individual relaxation and reaction processes to be examined
separately. They include pure nitrogen with and without electronic energy,
mixtures of molecular nitrogen with atomic nitrogen or oxygen, reacting gases
and five-species air. The results are compared with those from other CFD
solvers and direct simulation Monte Carlo (DSMC) codes, which represent the
gas through the motion and collisions of simulated particles. The present
work reproduces these heat-bath verification cases.

The remainder of this report is organised as follows.
Section~\ref{sec:lavoro} describes the software components, their coupling and
the model implementation. Section~\ref{sec:risultati} presents the
verification results for each case. Section~\ref{sec:discussione} interprets
the results and examines the remaining discrepancies, while
Section~\ref{sec:conclusioni} presents the conclusions and identifies
directions for future work.
